% Part: incompleteness
% Chapter: theories-computability
% Section: q-is-ce

\documentclass[../../../include/open-logic-section]{subfiles}

\begin{document}

\olfileid{inc}{tcp}{qce}

\olsection{$\Th{Q}$ تامة بين المجموعات \printtoken{S}{c.e.}}


\begin{thm}
$\Th{Q}$ !!{c.e.} وغير قابلة للقرار؛ وهي، فوق ذلك، تامة بين المجموعات
!!{c.e.}.
\end{thm}

\begin{proof}
أما كون $\Th{Q}$ !!{c.e.} فبيانه أنها مجموعة شفرات الجمل
${\OLId{latin-ordinary}{y}}$ التي يوجد برهان ${\OLId{latin-ordinary}{x}}$ عليها في $\Th{Q}$:
\[
{\OLId{latin-looped}{Q}} = \Setabs{{\OLId{latin-ordinary}{y}}}{\lexists[{\OLId{latin-ordinary}{x}}][\Prf[\Th{Q}]({\OLId{latin-ordinary}{x}},{\OLId{latin-ordinary}{y}})]}.
\]
والعلاقة $\Prf[\Th{Q}]({\OLId{latin-ordinary}{x}},{\OLId{latin-ordinary}{y}})$ قابلة للحساب، بل عودية بدائية؛ وقد
علمنا أن كل مجموعة ترد على هذه الصورة قابلة للتعداد حسابيًا.

وأما تمامها بين المجموعات القابلة للتعداد حسابيًا فمكافئ لكون
${\OLId{latin-looped}{K}}\leq_{\OLId{latin-ordinary}{m}} {\OLId{latin-looped}{Q}}$، حيث ${\OLId{latin-looped}{K}}=\Setabs{{\OLId{latin-ordinary}{x}}}{!A_{\OLId{latin-ordinary}{x}}({\OLId{latin-ordinary}{x}})\downarrow}$. فالمطلوب إذن
اختزال ${\OLId{latin-looped}{K}}$ إلى $\Th{Q}$. ومحمول كليني ${\OLId{latin-looped}{T}}({\OLId{latin-ordinary}{e}},{\OLId{latin-ordinary}{x}},{\OLId{latin-ordinary}{s}})$ عودي بدائي، ومن ثم
يقبل التمثيل في $\Th{Q}$؛ ولنسمّ الصيغة التي تمثله $!A_{\OLId{latin-looped}{T}}$. فلكل ${\OLId{latin-ordinary}{x}}$:
\begin{align*}
{\OLId{latin-ordinary}{x}} \in {\OLId{latin-looped}{K}} & \lif \lexists[{\OLId{latin-ordinary}{s}}][{\OLId{latin-looped}{T}}({\OLId{latin-ordinary}{x}},{\OLId{latin-ordinary}{x}},{\OLId{latin-ordinary}{s}})] \\
& \lif \lexists[{\OLId{latin-ordinary}{s}}][(\Th{Q} \Proves !A_{\OLId{latin-looped}{T}}(\num {\OLId{latin-ordinary}{x}},\num {\OLId{latin-ordinary}{x}}, \num {\OLId{latin-ordinary}{s}}))]
  \\
& \lif \Th{Q} \Proves \lexists[{\OLId{latin-ordinary}{s}}][!A_{\OLId{latin-looped}{T}}(\num {\OLId{latin-ordinary}{x}}, \num {\OLId{latin-ordinary}{x}}, {\OLId{latin-ordinary}{s}})].
\end{align*}
وأما العكس، فإذا كان
$\Th{Q}\Proves\lexists[{\OLId{latin-ordinary}{s}}][!A_{\OLId{latin-looped}{T}}(\num {\OLId{latin-ordinary}{x}},\num {\OLId{latin-ordinary}{x}},{\OLId{latin-ordinary}{s}})]$، لزم أن تصدق
الصيغة $!A_{\OLId{latin-looped}{T}}(\num {\OLId{latin-ordinary}{x}},\num {\OLId{latin-ordinary}{x}},\num {\OLId{latin-ordinary}{n}})$ عند عدد طبيعي ما ${\OLId{latin-ordinary}{n}}$.
فلو كانت ${\OLId{latin-looped}{T}}({\OLId{latin-ordinary}{x}},{\OLId{latin-ordinary}{x}},{\OLId{latin-ordinary}{n}})$ كاذبة، لأثبتت $\Th{Q}$ نفي تلك الصيغة،
$\lnot !A_{\OLId{latin-looped}{T}}(\num {\OLId{latin-ordinary}{x}},\num {\OLId{latin-ordinary}{x}},\num {\OLId{latin-ordinary}{n}})$، إذ $!A_{\OLId{latin-looped}{T}}$ ممثلة للعلاقة ${\OLId{latin-looped}{T}}$.
وحينئذ تكون $\Th{Q}$ قد أثبتت صيغة كاذبة، وهو محال. فثبت صدق
${\OLId{latin-looped}{T}}({\OLId{latin-ordinary}{x}},{\OLId{latin-ordinary}{x}},{\OLId{latin-ordinary}{n}})$، ولزم منه $!A_{\OLId{latin-ordinary}{x}}({\OLId{latin-ordinary}{x}})\downarrow$.

فالحاصل أن انتماء ${\OLId{latin-ordinary}{x}}$ إلى ${\OLId{latin-looped}{K}}$، لأي ${\OLId{latin-ordinary}{x}}$، يكافئ إثبات
$\Th{Q}$ للصيغة $\lexists[{\OLId{latin-ordinary}{s}}][!A_{\OLId{latin-looped}{T}}(\num {\OLId{latin-ordinary}{x}},\num {\OLId{latin-ordinary}{x}},{\OLId{latin-ordinary}{s}})]$.
وعليه فالدالة ${\OLId{latin-ordinary}{f}}$ التي تأخذ ${\OLId{latin-ordinary}{x}}$ وتعطي شفرة الجملة
$\lexists[{\OLId{latin-ordinary}{s}}][!A_{\OLId{latin-looped}{T}}(\num {\OLId{latin-ordinary}{x}},\num {\OLId{latin-ordinary}{x}},{\OLId{latin-ordinary}{s}})]$ هي اختزال ${\OLId{latin-looped}{K}}$ إلى
$\Th{Q}$ المطلوب.
\end{proof}

\end{document}
